% Generated draft: requires mathematical and source review.
\documentclass[11pt]{article}
\usepackage[margin=25mm]{geometry}
\usepackage{fontspec}
\setmainfont{Latin Modern Roman}
\usepackage{amsmath,amssymb,mathtools,unicode-math}
\setmathfont{Latin Modern Math}
\usepackage{enumitem,needspace,xurl,hyperref}
\setlength{\parindent}{0pt}
\setlength{\parskip}{5pt}
\setlength{\emergencystretch}{3em}
\newcommand{\N}{\mathbb{N}}
\newcommand{\Z}{\mathbb{Z}}
\newcommand{\Q}{\mathbb{Q}}
\newcommand{\R}{\mathbb{R}}
\newcommand{\C}{\mathbb{C}}
\newcommand{\F}{\mathbb{F}}
\newcommand{\A}{\mathbb{A}}
\newcommand{\PP}{\mathbb{P}}
\newcommand{\E}{\mathbb{E}}
\newcommand{\eps}{\varepsilon}
\newcommand{\dd}{\,\mathrm{d}}
\begin{document}
\section*{UnsolvedMath ID: 30006466}
\textbf{Source problem code:} catalogue entry
\subsection{Definitions and mathematical statement}
The supplied record combines several questions and contains an incomplete original excerpt. The following makes the cleaned statement precise under explicitly stated conventions; these conventions require source review.

Write \(\N_{>0}=\{1,2,3,\ldots\}\). Under the convention that the summands are positive integers, define the set of \(3\)-full numbers by
\[
B_3=\{m\in\N_{>0}: \text{for every prime }p\mid m,\ p^3\mid m\}.
\]
Thus every prime exponent in a member of \(B_3\) is at least \(3\), and \(1\in B_3\) by the vacuous condition. Define
\[
A=\{a+b+c:a,b,c\in B_3\}.
\]
Repeated summands are permitted in this reading. The supplied excerpt does not establish the conventions concerning positivity, the inclusion of \(1\), or repeated summands.

If ``positive density'' means positive natural density, the first question is whether there exists \(\delta\in(0,1]\) such that
\[
\lim_{X\to\infty}\frac{\#\bigl(A\cap[1,X]\bigr)}{\lfloor X\rfloor}=\delta.
\]
The source does not specify the density notion. For example, positive lower asymptotic density would instead mean
\[
\liminf_{X\to\infty}\frac{\#\bigl(A\cap[1,X]\bigr)}{\lfloor X\rfloor}>0.
\]
These formulations must not be identified without clarification. The second question is whether
\[
\exists N\in\N_{>0}\ \forall n\in\N_{>0},\qquad n\ge N\Longrightarrow n\in A.
\]
The original excerpt includes the informal expectation that the answer to this second question is negative and that the density of \(A\) is probably less than \(1\). These are supplied expectations, not established conclusions. Its opening question, ``Is this easier to prove for the sum of three 3-full numbers?'', has no supplied antecedent; the positive-density question is explicit only in the cleaned statement.

Separately, for \(s\in\C\), define
\[
f(s)=1+2^{-s}+3^{-s}+5^{-s},\qquad n^{-s}=\exp(-s\log n),
\]
where \(\log n\) is the real logarithm of the positive integer \(n\). The zero question is
\[
\exists t\in\R\quad f(1+it)=0,
\]
where \(i^2=-1\). The supplied excerpt attributes this question to Erd\H{o}s and Ingham in connection with Tauberian theorems and records the informal expectation that no such \(t\) exists.

The original excerpt also contains adjacent questions about the number of palindromic squares in \([1,X]\), and about whether infinitely many non-palindromic squares remain squares after reversal of their digits, with examples \(144\) and \(169\) and the identifier OEIS A033294. It asserts that infinitely many palindromic squares exist. An editor's note says that Keith [9] gave four infinite families and asked whether infinitely many palindromic squares lie outside those families. The numeral base, reversal conventions, the families, and the bibliographic identity of [9] are not supplied. The nearby attribution ``Proposed by Julia Brandes'' has unclear scope. These adjacent questions are retained as unresolved source context; the record does not establish that they belong to the intended target of this permanent ID.
\subsection{Short English statement}
Under explicit positive-integer and natural-density conventions, ask whether sums of three \(3\)-full numbers have positive density and whether every sufficiently large positive integer is such a sum. Separately, ask whether \(1+2^{-s}+3^{-s}+5^{-s}\) has a zero on \(\operatorname{Re}s=1\). The original excerpt has missing context and includes additional digit-reversal questions whose scope is unresolved. The supplied open label and informal expectations require independent review.
\subsection{Sources}
\begin{itemize}
\item Supplied local UnsolvedMath registry record, permanent ID \(30006466\), exact source problem code \texttt{catalogue entry}, titled ``Three-Full Sums and the Erd\H{o}s--Ingham Dirichlet Zero Problem''. The supplied collection is ``Oberwolfach Reports Open Problems'', with proposed year \(2025\).
\item Supplied report citation: \emph{Analytic Number Theory (2026)}, Oberwolfach Reports, DOI \texttt{10.4171/owr/2025/51}, pp.\ 2701--2764. Source URL and supplied literature-source link, labeled ``Oberwolfach report 51/2025'': \url{https://doi.org/10.4171/owr/2025/51}. The differing year indications are retained as supplied.
\item Supplied literature assessment, checked exactly \texttt{2026-08-22}: the record labels the questions open and reports that its curation found no authoritative proof resolving either the three-full-sum questions or the Erd\H{o}s--Ingham zero question. This is supplied dated curation, not an independently verified literature assessment; time-sensitive status and cited claims require independent verification.
\item The original excerpt attributes the zero question to Erd\H{o}s and Ingham in connection with Tauberian theorems. It also includes ``Proposed by Julia Brandes'', whose scope is unclear, and an editor's reference to Keith [9] concerning four infinite families of palindromic squares. No full citation for [9] is supplied. OEIS A033294 is supplied as an identifier for the digit-reversal question, without an accompanying link.
\end{itemize}
\end{document}
